\documentclass[../main.tex]{subfiles}

\begin{document}

\ifSubfilesClassLoaded{

    \tableofcontents
    
}{}

\chapter{ Smooth Maps }

\section{Smooth Functions and Smooth Maps}

\subsection{Smooth Maps Between Manifolds}

\begin{definition}{Smooth Maps Between Manifolds}{}
    Let $M, N$ be smooth manifolds, and let $F: M \to N$ be any map. We say that $F$ is a \dfntxt{smooth map} if for every $p \in M$, there exist smooth charts $(U, \varphi)$ containing $p$ and $(V, \psi)$ containing $F(p)$ such that $F(U) \subseteq V$ and the composite map $\psi \circ F \circ \varphi^{-1}$ is smooth from $\varphi(U)$ to $\psi(V)$ .
\end{definition}

\begin{remark}
    \begin{itemize}
        \item By taking \(  N= V= \mathbb{R} ^{k}  \) and \(  \psi =\mathrm{Id}:\mathbb{R} ^{k}\to \mathbb{R} ^{k}  \), we get the definition of smoothness of real-valued or vector-valed functions.  
    \end{itemize}
    
\end{remark}



\section{Partitions of Unity}
\begin{definition}{Support}{}
    If $f$ is any real-valued or vector-valued function on a topological space $M$, the \dfntxt{support of f}, denoted by $\operatorname{supp} f$, is the closure of the set of points where $f$ is nonzero:
\[
\operatorname{supp} f = \overline{\{p \in M : f(p) \neq 0\}}.
\]\begin{itemize}
    \item If $\operatorname{supp} f$ is contained in some set $U \subseteq M$, we say that $f$ is \dfntxt{supported in U}. 
    \item 
A function $f$ is said to be \dfntxt{compactly supported} if $\operatorname{supp} f$ is a compact set, denoted by \(  f \in C_{0}^{\infty}\left( M \right)   \) . 
\end{itemize}
\end{definition}

\begin{remark}
    Clearly, every function on a compact space is compactly supported.
\end{remark}

\begin{lemma}{}{}
The function $f: \mathbb{R} \to \mathbb{R}$ defined by
\[
f(t) = \begin{cases}
e^{-1/t}, & t > 0, \\
0, & t \le 0,
\end{cases}
\]
is smooth.
\end{lemma}
\begin{lemma}{}{}
Given any real numbers $r_1$ and $r_2$ such that $r_1 < r_2$, there exists a smooth function $h: \mathbb{R} \to \mathbb{R}$ such that $h(t) \equiv 1$ for $t \le r_1$, $0 < h(t) < 1$ for $r_1 < t < r_2$, and $h(t) \equiv 0$ for $t \ge r_2$.
\end{lemma}
\begin{proof}
Let $f$ be the function of the previous lemma, and set
\[
h(t) = \frac{f(r_2 - t)}{f(r_2 - t) + f(t - r_1)}.
\]Note that the denominator is positive for all $t$, because at least one of the expressions $r_2 - t$ and $t - r_1$ is always positive. The desired properties of $h$ follow easily from those of $f$.
\end{proof}

\begin{lemma}{Bump Function on \(  \mathbb{R} ^{n}  \) }{9-17-1}
Given any positive real numbers $r_1 < r_2$, there is a smooth function $H : \mathbb{R}^n \to \mathbb{R}$ such that $H \equiv 1$ on $\bar{B}_{r_1}(0)$, $0 < H(x) < 1$ for all $x \in B_{r_2}(0) \setminus \bar{B}_{r_1}(0)$, and $H \equiv 0$ on $\mathbb{R}^n \setminus B_{r_2}(0)$.
\end{lemma}

\begin{proof}
Just set $H(x) = h(|x|)$, where $h$ is the function of the preceding lemma. Clearly, $H$ is smooth on $\mathbb{R}^n \setminus \{0\}$, because it is a composition of smooth functions there. Since it is identically equal to 1 on $B_{r_1}(0)$, it is smooth there too. 
\end{proof}

\begin{theorem}{Existence of Bump Function on Manifold}{}
    Let \(  M  \) be a smooth manifold, \(  K\subseteq M  \) a compact set. \(  U\subseteq M  \) is a open set such that \(  K\subseteq U  \). Then there exists \(   \varphi  \in C_{0}^{\infty}\left( M \right)   \), such that \(  \operatorname{supp} \varphi \subseteq U  \), \(  0\le  \varphi \le 1  \) and \(   \varphi \equiv 1  \) on \(  K  \).           
\end{theorem}

\begin{proof}
    For each \(  x \in K  \), take a chart \(  \left(  \varphi _{x}, U_{x},V_{x} \right)   \) centering \(  x  \). We may assume that \(  V_{x} \subseteq B_{3}\left( 0 \right)   \) and \(  U_{x}\subseteq U  \)  by contracting  when necessary . Let \(  \tilde{U_{x}}  \coloneqq  \varphi _{x}^{-1} \left( B_{1}\left( 0 \right)  \right) \). We define \[
    f_{x}\left( p \right)\coloneqq \begin{cases} H\circ  \varphi _{x}\left( p \right),&p\in U_{x}\\ 
     0,& p \not \in U_{x}  \end{cases}  
    \]   Where \(  H  \) is defined in Lemma \ref{lem:9-17-1}  and we take \(  r_1= 1, r_2= 2  \).   It is straightforward to verify that \(  f_{x}\equiv 1  \) on \(  \tilde{U_{x}}  \) and \(  f_{x}\equiv 0  \) on \(  U_{x}^{c}  \), \(  \operatorname{supp}f_{x}=  \varphi _{x}^{-1} \left( \overline{B_2\left( 0 \right) } \right)  \subseteq U_{x} \). Since the homeomorphism keeps the closure \(   \varphi _{x}^{-1} \left( \overline{B_{2}\left( 0 \right) } \right)= \overline{ \varphi _{x}^{-1} \left( B_{2}\left( 0 \right)  \right) }   \). 
  
    Now \(  K  \) has a open cover \(  \left\{ \tilde{U}_{x} \right\}  \). Since \(  K  \) is compact, there are finitely many \(  \tilde{U}_{x_{i}}  \) such that \[
    \bigcup _{i= 1}^{N} \tilde{U}_{x_{i}}\supseteq K
    \]    We define \[
    \psi \coloneqq \sum _{i= 1}^{N}f_{x_{i}}
    \]Then \(  \psi \ge 1 \) on \(  K  \), \(  \operatorname{supp}\psi \subseteq U  \).    

    Let $\rho: \mathbb{R} \to [0, 1]$ be a smooth function such that
$$\rho(t) = \begin{cases} 
0, & t \le 1/2 \\
1, & t \ge 1
\end{cases}$$
and $\rho$ is monotonically increasing on $(1/2, 1)$. 
Now, we define our final function $\varphi: M \to \mathbb{R}$ as the composition
$$\varphi \coloneqq \rho \circ \psi $$
\end{proof}

\begin{definition}{Partitions of Unity}{}
    Suppose $M$ is a topological space, and let $\mathcal{X} = (X_\alpha)_{\alpha \in A}$ be an arbitrary open cover of $M$, indexed by a set $A$. A \dfntxt{partition of unity subordinate to $\mathcal{X}$} is an indexed family $(\psi_\alpha)_{\alpha \in A}$ of continuous functions $\psi_\alpha: M \to \mathbb{R}$ with the following properties:
\begin{enumerate}
\item $0 \le \psi_\alpha(x) \le 1$ for all $\alpha \in A$ and all $x \in M$.
\item $\text{supp } \psi_\alpha \subseteq X_\alpha$ for each $\alpha \in A$.
\item The family of supports $(\text{supp } \psi_\alpha)_{\alpha \in A}$ is locally finite, meaning that every point has a neighborhood that intersects $\text{supp } \psi_\alpha$ for only finitely many values of $\alpha$.
\item $\sum_{\alpha \in A} \psi_\alpha(x) = 1$ for all $x \in M$.
\end{enumerate}

If $M$ is a smooth manifold with or without boundary, a \dfntxt{smooth partition of unity} is one for which each of the functions $\psi_\alpha$ is smooth.
\end{definition}

\begin{remark}
    \begin{itemize}
    \item 
Because of the local finiteness condition (3), the sum in (4) actually has only finitely many nonzero terms in a neighborhood of each point, so there is no issue of convergence. 
\item For a manifold \(  M  \),   \(  M  \) is Lindel\"of Space, we can always assmue that \(  \left( X_{\alpha } \right) _{\alpha \in A}  \) is a countable colletction, thus the \(  \left( \psi _{\alpha } \right)_{\alpha \in A}   \) is as well .  
    \end{itemize}
    
\end{remark}

\begin{theorem}{Existence of Partitions of Unity}{}
Suppose $M$ is a smooth manifold with or without boundary, and $\mathcal{X} = (X_{\alpha})_{\alpha \in A}$ is any indexed open cover of $M$. Then there exists a smooth partition of unity subordinate to $\mathcal{X}$.
\end{theorem}
\subsection*{Applications}

\begin{proposition}{Bump function Support on a Closed Set}{}
    Let \(  M  \) be a manifold, \(  A\subseteq M  \) be a closed set, \(  U\supseteq A  \) is an open set. Then there exists \(   \varphi  \in C_{0}^{\infty}\left( M \right)   \) such that \(  0\le  \varphi \le 1  \), and  \(   \varphi \equiv 1  \) on \(  A  \)       .
\end{proposition}

\begin{proof}
    \(  \left\{ U, M\setminus A \right\}  \) is a open cover of \(  M  \). Then there exists a partition  of unity \(  \left\{ \rho _1 ,\rho _2  \right\}  \) subordinate to it.  Then \(  \rho _1   \) is just the function we need. 
\end{proof}


\begin{theorem}{Stone-Weierstrass}{}
    Let \(  X  \) be a compact Hausdorff Space. \(  \mathcal{A}\subseteq C\left( X,\mathbb{R}  \right)   \) is a subalgebra . Then \(  \mathcal{A}  \) is dense in \(  C\left( X,\mathbb{R}  \right)   \), if and only if it satisfies  
    \begin{enumerate}
        \item Nowhere vanishing: For each \(  x \in X  \), there exists \(  g \in \mathcal{A}  \) such that \(  g\left( x \right)\neq 0   \).
        \item Separates points: For distinct \(  x,y \in X  \), there exists \(  f\in \mathcal{A}  \) such that \(  f\left( x \right)\neq g\left( x \right)    \).      
    \end{enumerate}
    
\end{theorem}
\begin{theorem}{Whitney}{}
    Let \(  M  \) be a smooth manifold. For each continuous function \(  g\in :M\to \mathbb{R}   \), \(   \delta :M\to \mathbb{R} _{+ }  \)   , there exists a smooth function \(  f :M\to \mathbb{R}  \) , such that   \[
    \left| f\left( p \right)-g\left( p \right)   \right|<  \delta \left( p \right),\quad \forall p\in M  
    \]
\end{theorem}

\begin{theorem}{Whitney Extension}{}
    Let \(  M  \) be a smooth manifold.  \(  A\subseteq M  \) is a closed set. For continuous \(  g: M\to \mathbb{R}   \), \(   \delta :M\to \mathbb{R} _{+ }  \)  , where \(  g  \) is smooth on \(  A  \). Then there exists smooth function \(  f:M\to \mathbb{R}   \), such that
    \begin{enumerate}
        \item \[
        f\left( p \right)= g\left( p \right),\quad \forall p\in A  
        \]
        \item \[
        \left| f\left( p \right)-g\left( p \right)   \right|<  \delta \left( p \right),\quad \forall p \in M  
        \]
    \end{enumerate}
     
\end{theorem}
\end{document}